% New proof for the combined AX11/Q/G/H interface.  The inherited numerical
% calibration and generic source/tape lemmas are identified explicitly below.
% Requires amsmath, amssymb; this file is a section fragment, not a standalone preamble.
\section{PA194AX11QGH: eleven entries, direct binder, and cached halting}
\label{sec:ax11qgh}

The new table combines three independently observable command changes:
AX saves eleven entries and consumes a fourth operand; its saved
generalization uses its explicit first operand as binder; and Shift
tests only the cached formula.  Its sixth logical head has the existing
Q orientation.  The abbreviation \textsf{AX11QGH} specifies all four
features.  Earlier names such as \textsf{AX11G} already included Q; we
make that inherited choice explicit rather than count it as a new idea.

The table has 194 working states, start 0, halt 194, 388 rows, and a
two-way binary blank tape.  Its SHA-256 is
\begin{quote}\small\texttt{513e9c8b28e0dd41e6003bb8d82c80ff3dac5fb3c5274cd2f709a3bba8c3e9e1}.
\end{quote}
The ordinary mathematical theorem below reuses the written weak-base
calibration in the incoming \textsf{pa194\_ax11h\_research\_2026-11-14}
packet at ProveIt commit
\texttt{1e1a41f799d4658b829cdd0eabae3775a0c8c4a0}.
The new executable certificates establish the candidate-specific
finite and symbolic identities.  They are not an end-to-end Lean theorem.

\subsection{Numerical calculus and the complete store transformation}

Let
\[
 p(a,b)=\frac{(a+b)(a+b+1)}2+b,\quad [a;w]=p(a,w),
\]
and use total Cantor projections.  Thus
$\operatorname{hd}[a;w]=a$, $\operatorname{tl}[a;w]=w$, and
$[0;0]=\operatorname{hd}(0)=\operatorname{tl}(0)=0$.
Put
\[
 I(A,B)=4p(A,B),\quad M(a,b)=4p(a,b)+1,\quad
 U(a,B)=4p(B,a)+2,\quad E(a,b)=4p(a,b)+3.
\]
Write $A\to B=I(A,B)$, $N(A)=I(A,1)$,
$A\wedge B=N(I(A,N(B)))$, and $\exists_a A=N(U(a,N(A)))$.
Zero is truth, and one is the contradiction code.  The retained heads are
\begin{align*}
 h_1(D,H)&=(\exists_D U(D,H))\to H,\\
 h_2(D,A,H)&=U(D,U(A,H))\to U(A,U(D,H)),\\
 h_3(D,A,H)&=M(H,D)\to(E(H,A)\to M(A,D)),\\
 h_4(D,A,H)&=E(H,D)\to(M(A,H)\to M(A,D)),\\
 h_5(D,H)&=\exists_D E(H,D),\\
 h_6^Q(D,A,H)&=U(D,I(A,H))\to(U(D,A)\to U(D,H)),\\
 \Gamma(A)&=A\to U(A,A),\\
 \rho(D,A,H,u)&=((H\to u)\to D)\to((D\to H)\to(A\to H)),
 \qquad u\le D.
\end{align*}
For the guarded gate define
\begin{align*}
 V(a,b,c)&=M(a,b)\to U(c,M(a,b)), & P(a,b,c)&=c\to M(c,b),\\
 S(a,b,c)&=U\bigl(b,(M(a,b)\to P(a,b,c))\wedge(M(c,b)\to c)\bigr),\\
 Q(a,b,c)&=U(c,S(a,b,c)\to N(P(a,b,c))),\\
 g(a,b,c)&=\begin{cases}N(V(a,b,c)\to Q(a,b,c)),&c\ne a,b,\\
 0,&\text{otherwise.}\end{cases}
\end{align*}
These are the same axiom families and the same MP and theorem-generalization
rules as the incoming QAX11H calculus $\mathcal C_4^Q$.
The instance $\rho(0,0,0,0)=0$ is a fixed proof of truth.

The command store is $(t,w,d,h)$: visible theorem, numerical stack,
offset, and cached theorem.  For $D=d+x$ and $H=h+y$, AX has exactly
the effect
\begin{equation}
 A_{c,x,y,z}(t,w,d,h)
   =(0,[q_0,q_1,\ldots,q_9,t;w],z,0),
 \tag{QGH-A}
\end{equation}
where
\[
\begin{array}{c|l@{\qquad}c|l}
0&g(H,D,c)&1&\Gamma(c)\\
2&h_6^Q(D,c,H)&3&h_5(D,H)\\
4&h_4(D,c,H)&5&h_3(D,c,H)\\
6&h_2(D,c,H)&7&h_1(D,H)\\
8&\rho(D,c,H,d)&9&U(c,h).
\end{array}
\]
In particular $q_9$ uses the incoming theorem $h$, not $h+y$;
its binder is $c$, not $d$.  Both distinctions are necessary.
The fourth operand is a genuine field: the added push leaves
$t=d=h=0$ before the full MP parameter reader.  The subsequent comparison
is $I(0,z)=0$.  If $z=0$ it succeeds and returns zero; if $z>0$
it fails and returns zero.  Both branches leave offset $z$ and preserve
the extended stack.  Thus (QGH-A) holds for arbitrary $z$.

All $q_0,\ldots,q_8$ have residue zero modulo four and $q_9$ has
residue two.  No generated entry is one.  The other commands are
\[
 S(t,w,d,h)=\begin{cases}\mathsf{halt},&h=1,\\
 (h,\operatorname{tl}w,d,\operatorname{hd}w),&h\ne1,
 \end{cases}
\]
and, writing $b=d+x$,
\[
 \operatorname{MP}_x(t,w,d,h)=
 \begin{cases}(b,w,b,0),&I(t,b)=h,\\(0,w,b,0),&I(t,b)\ne h.
 \end{cases}
\]
MP has no dummy push.  Only Shift can halt.

\subsection{Exact local words, including the zero aliases}

Words execute left to right.  Set $A_0=A_{0,0,0,0}$ and
\[
 \operatorname{Push}=A_0S^{10}\operatorname{MP}_0,\quad
 \operatorname{Pop}=A_0S^{12},\quad
 \operatorname{Discard}=\operatorname{MP}_0.
\]
Direct projection calculations give
\begin{align*}
 A_{c,x,y,0}S^{10}(t,w,d,h)&=(q_8,[t;w],0,U(c,h)),\\
 A_{c,x,y,0}S^{11}(t,w,d,h)&=(U(c,h),w,0,t),\\
 \operatorname{Push}(t,w,d,h)&=(0,[t;w],0,0),\\
 \operatorname{Pop}(t,[a;w],d,h)&=
   \begin{cases}\mathsf{halt},&t=1,\\(t,w,0,a),&t\ne1,
   \end{cases}\\
 \operatorname{Discard}(t,w,0,h)&=(0,w,0,0),\\
 \operatorname{MP}_i(0,w,0,0)&=(0,w,i,0).
\end{align*}
The first eleven Shifts after AX test zero and generated entries, never
the saved top.  The twelfth tests that top.  Push's final comparison
fails by the distinct implication/universal residues.  These arguments
use the total numerical projections; there is no assumption that a
numerical push increases a free-list length.

For $0\le j\le8$, the word
\[
 \operatorname{Discard};A_{c,D,H,0};S^{j+2};
                   \operatorname{Pop}^{9-j}
\]
returns $(q_j,w,0,0)$ from $(t,w,0,0)$, with the parameters displayed
in the table.  The selected entry is not one, so its Pops are safe;
the last removed entry is the saved zero.  To select a restricted
$\rho(D,c,H,u)$ with $u\le D$, use
\[
 \operatorname{Discard};\operatorname{MP}_u;
 A_{c,D\mathbin{\dot{-}}u,H,0};S^{10};\operatorname{Pop}.
\]
The inequality gives $u+(D\mathbin{\dot{-}}u)=D$.

The new generalization word is
\begin{equation}
 \operatorname{Gen}_i=\operatorname{Push};S;
                   A_{i,0,0,0};S^{11}.
 \tag{QGH-Gen}
\end{equation}
It sends $(a,w,0,0)$ to $(U(i,a),w,0,0)$ for every $a,w,i$.
After Push and the first Shift, the store is $(0,w,0,a)$.
AX consumes the cache as a theorem before another Shift can test it.
Eleven Shifts then select $U(i,a)$ and the saved zero becomes the cache.
The claim includes $a=1$ and $i=0$.  The incoming QAX11H compiler has
an additional offset-setting command; (QGH-Gen) removes it because
the binder now comes directly from AX's operand.

The MP suffix
\[
 S;S;\operatorname{MP}_b;\operatorname{Pop}
\]
on $(0,[a,I(a,b),0;w],0,0)$ halts if $a=1$ or $b=1$, and otherwise
returns $(b,w,0,0)$.  The second Shift tests $a$.  After a successful
MP, Pop removes the reserved zero and tests $b$.  Omitting that
reserved zero would consume an arbitrary part of the underlying stack.

\subsection{Proof compilation in quantifier-free induction}

Compile numerical derivations with the preceding selectors and
generalization word.  At an MP node with minor proof $e_0$ of $a$
and major proof $e_1$ of $I(a,b)$, use
\[
 \operatorname{Discard};\operatorname{Push};W_{e_1};
 \operatorname{Push};W_{e_0};\operatorname{Push};
 S;S;\operatorname{MP}_b;\operatorname{Pop}.
\]
Its first Discard and Push reserve a zero.  The exact invariant is:
from every $(t,w,0,0)$, $W_e$ either halts or returns
$(\operatorname{concl}e,w,0,0)$.  The local identities prove the
axiom cases and each composition step.

Here is an explicit induction-strength justification.  Encode a valid
derivation as $N$ lines whose premises have earlier indices.  Form the
compiled words $W_j$ by primitive recursion on $j$, storing earlier
words.  Let $L$ be their maximum length and $K$ their maximum command
code.  Totalize command execution with an absorbing halt flag and define
\[
 \Phi_K(s)=1+\max\left(s,
       \max_{c\le K,\ u\le s}\operatorname{next}(c,u)\right).
\]
This is a monotone primitive-recursive majorant.  Fix an arbitrary
normalized input code $s_0$ and let
$B(j)=\Phi_K^{L(N-1-j)}(s_0)$ for $j<N$.
For $i<j$, $\Phi_K^L(B(j))\le B(i)$.
At stage $j$ retain the finite conjunction asserting the required
execution equality for every earlier line $i<j$ and every normalized
store code at most $B(i)$.  Its characteristic function is primitive
recursive, so this is an allowed quantifier-free induction invariant.
Every recursive premise call occurs after a prefix of at most $L$
commands and falls below the earlier bound.  The local word identities
therefore establish the successor step.  At $j=N-1$ the bound is
$s_0$; universal generalization gives the all-store statement.
This does not appeal to induction on an unbounded assertion about all
stacks.

If $e$ is a refutation, then $A_0;W_e;\operatorname{Pop}$ halts
from every incoming store.  The first AX clears the two caches.
If $W_e$ returns top one rather than halting earlier, Pop detects it.
A final lone Shift is insufficient because it tests the zero cache.

\subsection{Scheduler completeness and soundness}

For a word let $e(\epsilon)=0$ and $e(c::w)=1+b_c(e(w))$, where
\[
 b_S(z)=p(0,z),\quad b_{\operatorname{MP}_x}(z)=p(1,p(x,z)),
\]
\[
 b_{A_{c,x,y,u}}(z)=p(c+2,p(y,p(x,p(u,z)))).
\]
These are four-field AX codes.  Each decoder returns a nested component
of its seed.  With scheduler state $(n,\ell,t,w,d,h)$, the seed is
$n$ and the fresh counter becomes $n+1$ when $\ell=0$; otherwise
the seed is $\ell-1$ and the counter is unchanged.  The initial
state is $(2,0,0,0,0,0)$.
Every positive continuation decreases.  A fresh candidate $n$ drains
in at most $n+1$ commands unless it halts earlier.
If $z=e(A_0;W_e;\operatorname{Pop})$, the fresh candidate
\[
 N_e=p(2,p(0,p(0,p(0,z))))\ge3
\]
executes an AX prefix followed by the all-store halting suffix.
The bound $(N_e+2)^2$ suffices to reach and drain it from the initial
candidate two.  The carried stack is never assumed to reset.
Decoder inversion, decreasing continuation, and this bound have
primitive-recursive Boolean formulations and quantifier-free proofs.

Conversely maintain a finite numerical-calculus proof list annotating
$t$, $h$, and every canonical stack entry.  Truth annotates virtual zero
tails.  AX adds its fixed axiom instances, its saved top, and the theorem
$U(c,h)$ obtained by generalizing the incoming proof of $h$.
The restricted head has its valid bound $d\le d+x$; the failed gate
uses truth.  The offset is not required to be a theorem.
Successful MP combines the two annotated proofs; failure adds truth.
Shift only moves annotations and exposes a proof of one when it halts.
Canonical stack decoding is bounded by $w$ since $\operatorname{tl}w<w$
for $w>0$.  All annotation operations and validity checks are primitive
recursive, so quantifier-free induction on the trace proves soundness.
Consequently iPRA proves that this abstract scheduler halts exactly
when $\mathcal C_4^Q$ has a refutation.

\subsection{The exact source and physical bridges}

The frozen AX11G grammar has 67 reader nodes and already has the final
gate push.  Change only the five H Shift rows and the source node
\[
 \texttt{node170}=\texttt{((node36,(node167,node145)),node57)}.
\]
The rows are $(38,0):0R11$, $(38,1):1R122$,
$(214,0):0R312$, $(214,1):1R20$, and $(312,0):0R188$.
The seven retained disjoint contextual folds give $134+67-7=194$.
Every one of their 28 two-cell cases is checked on the final simultaneous
table.  Their exterior tape is arbitrary.

The builder also populates the accepted QAX11H source association with
the one new macro substitution \texttt{copyD} to \texttt{copyA} at
zero-based token 1.  Both complete resolved primitive source graphs
agree: all 3,671 instructions and all successor indices coincide.
The arithmetic macro itself is the same Copy template at another
register; the generic Copy proof is uniform in that register.
The free-algebra checker follows the entire 190-token emitter and verifies
every saved formula, the gate, the guard suffix and the genuine fourth
operand.  Thus the source-command proof changes only $q_9$ and its
shorter generalization word; all source-loop clocks are the same
primitive-recursive templates with the changed register value as input.

The independent literal checker reconstructs four unary-list sweep
lemmas from the actual rows and checks all 4,778 source outcomes for
arbitrary nine natural-number registers, both decrement branches,
arbitrary nonempty finite spare lists and arbitrary extra header gap.
It checks the entire blank-start tape after 57,922 transitions.
These are symbolic instruction contracts, not a long sampled run.
The separate frame-closure audit checks all 4,777 nonhalting outcomes
against the independently resolved source successor graph.  Every
48-bit padded target address agrees, every overflow is zero, and every
13-component affine head-motion vector is zero.  Every duration vector
has nonnegative coefficients and constant at least three; the unique
halt outcome has duration exactly three.  The maximum source path
depth is 34.  Thus the local endpoint and positivity conditions used
by the composition theorem are checked explicitly for this table.
The generic physical bridge of the incoming AX11H proof applies to
these identical frame conventions and newly replayed contracts:
iterate the explicit source frame and sum its positive primitive-
recursive durations.  A physical first halt at time $T$ is reflected
by at most $T+1$ nonhalting source steps, since otherwise their positive
segments continue beyond the first halt.  This is bounded induction,
including nonstandard model times, rather than standard-run reasoning.

The written AX11H/QAX11H packet supplies the separate
Foundation-PA proof adapter, numerical-calculus calibration in
$I\Sigma_1$, and the transfer to iPRA for the negative sentences.
No axiom family is changed here, and the previous subsections supply
both new execution bridges.  Composing them with that stated written
calibration therefore proves
\[
 \mathrm{iPRA}\ \vdash\quad
 \operatorname{Con}(\mathrm{PA})\ \longleftrightarrow\
 \neg\operatorname{Halt}_{\mathrm{blank}}(\mathrm{PA194AX11QGH}).
\]
This is a written mathematical proof with inherited, identified
calibration dependencies and new executable bindings.  The incomplete
end-to-end Lean development is not claimed to have been completed.
